% Generated draft: requires mathematical and source review.
\documentclass[11pt]{article}
\usepackage[margin=25mm]{geometry}
\usepackage{fontspec}
\setmainfont{Latin Modern Roman}
\usepackage{amsmath,amssymb,mathtools,unicode-math}
\setmathfont{Latin Modern Math}
\usepackage{enumitem,needspace,xurl,hyperref}
\setlength{\parindent}{0pt}
\setlength{\parskip}{5pt}
\setlength{\emergencystretch}{3em}
\newcommand{\N}{\mathbb{N}}
\newcommand{\Z}{\mathbb{Z}}
\newcommand{\Q}{\mathbb{Q}}
\newcommand{\R}{\mathbb{R}}
\newcommand{\C}{\mathbb{C}}
\newcommand{\F}{\mathbb{F}}
\newcommand{\A}{\mathbb{A}}
\newcommand{\PP}{\mathbb{P}}
\newcommand{\E}{\mathbb{E}}
\newcommand{\eps}{\varepsilon}
\newcommand{\dd}{\,\mathrm{d}}
\begin{document}
\section*{UnsolvedMath ID: 800004}
\textbf{Source problem code:} catalogue entry
\subsection{Definitions and mathematical statement}
For a positive integer \(n\), let \([n]=\{1,\ldots,n\}\) and let
\[
T(n)=\{f:[n]\to[n]\}
\]
be the full transformation semigroup, with multiplication given by composition, \((fg)(a)=f(g(a))\). The positive-integer range is the usual convention adopted here; the supplied source does not explicitly specify the range of \(n\), including whether an empty underlying set is intended.

Let \(X=\{x_1,x_2,\ldots\}\), and let \(X^+\) be the free semigroup of finite nonempty words over \(X\), with concatenation as multiplication. Each assignment \(\phi:X\to T(n)\) extends uniquely to a semigroup homomorphism \(\widehat\phi:X^+\to T(n)\). A semigroup identity \(u\approx v\), where \(u,v\in X^+\), holds in \(T(n)\) if
\[
\widehat\phi(u)=\widehat\phi(v)
\qquad\text{for every assignment }\phi:X\to T(n).
\]
Thus its equational theory is
\[
\operatorname{Id}(T(n))=
\bigl\{(u,v)\in X^+\times X^+:
\widehat\phi(u)=\widehat\phi(v)
\text{ for every }\phi:X\to T(n)\bigr\}.
\]
\textbf{Problem (E. S. Lyapin).} Describe \(\operatorname{Id}(T(n))\) for an \(n\)-element underlying set, with the qualification on the range of \(n\) stated above. The source asks for a description of the identities; it does not specify a required form of description or ask explicitly for a finite identity basis. Nonempty words are used because the source specifies semigroup identities.

The supplied source reports that some classes of these identities were described in works labeled [PSSSV] and [Ma3]. Their bibliographic details and the classes covered are absent from the supplied record, so this reported partial progress cannot be made more precise here.
\subsection{Short English statement}
Describe exactly which equations between nonempty words hold under every substitution of transformations of an \(n\)-element set. The problem is attributed to E. S. Lyapin. The supplied registry labels it open and reports partial descriptions in [PSSSV] and [Ma3]; these are supplied source assessments, not independently verified conclusions.
\subsection{Sources}
\begin{itemize}
\item Local UnsolvedMath registry, permanent ID \(800004\), source problem code \texttt{catalogue entry}, title \emph{The Lyapin's Notebook --- Problem 4}; supplied complete source record, upstream revision \texttt{d371ee917ac9647fbb06cd6fc529deec6f3dda82}. The record was created and updated on 2026-10-16 and carries the status \texttt{open}.
\item Bershadsky, Kublanovsky, Mashevitzky, \emph{The Lyapin's Notebook: Unsolved Problems in Semigroup Theory} (2026), Problem 4, \url{https://arxiv.org/abs/2604.04763}. The supplied provenance records access on 2026-10-16 and extraction by \texttt{arxiv\_latex}; this access information is supplied provenance, not a claim of independent consultation. The statement attributes the problem to E. S. Lyapin; the registry's proposer field names Bershadsky, Kublanovsky, Mashevitzky and gives the year 2026.
\item [PSSSV]: reference label supplied for work reportedly describing some classes of identities of \(T(n)\). No bibliographic details or separate link are supplied.
\item [Ma3]: reference label supplied for work reportedly describing some classes of identities of \(T(n)\). No bibliographic details or separate link are supplied.
\end{itemize}
\end{document}
